\documentclass[11pt,a4paper]{article}

\usepackage{fontspec}
\setmainfont{Libertinus Serif}
\usepackage[greek.ancient,spanish]{babel}
\babelfont[greek]{rm}{Libertinus Serif}
\usepackage[a4paper,margin=3.2cm]{geometry}
\usepackage{microtype}
\usepackage[hidelinks]{hyperref}

% Griego politónico escrito directamente en UTF-8:
% \gr{...} para palabras sueltas; bloques con \foreignlanguage{greek}{...}
\newcommand{\gr}[1]{\foreignlanguage{greek}{#1}}
\newcommand{\stanzabreak}{\par\vspace{0.6\baselineskip}}

\setlength{\parindent}{1.25em}
\setlength{\parskip}{0.35em}
\raggedbottom

\usepackage{enumitem}

\hypersetup{
  pdftitle={Leer la Odisea en tiempos iletrados: resumen de capítulos},
  pdfauthor={Juan José Gómez Cadenas}
}

\title{Leer la \emph{Odisea} en tiempos iletrados\\[0.4em]
\large Resumen de capítulos}
\author{Juan José Gómez Cadenas}
\date{}

\begin{document}

\maketitle

El libro propone una lectura personal de la \emph{Odisea} a partir de la
traducción de Juan Manuel Rodríguez Tobal. La película de Christopher Nolan
sirve como ejemplo de las dificultades del cine comercial para conservar la
complejidad de Homero. El análisis se entrelaza con recuerdos del autor y con
su poema \emph{Abandonando Ítaca}.

\begin{enumerate}[label=\Roman*.,leftmargin=3.2em,itemsep=0.7em]

\item \textbf{De libros y películas.}
La pérdida de hábitos lectores hace que el cine sea para muchos el principal
acceso a los clásicos. La coincidencia entre la película de Nolan y la
traducción de Tobal da origen al libro.

\item \textbf{El cine teme a Homero.}
El recuerdo de la primera lectura de Homero conduce al problema central. El
cine de masas tiende a sustituir la ambigüedad de los héroes por personajes
moralmente reconocibles.

\item \textbf{Una osada \emph{Odisea}.}
La traducción de Tobal recupera la música del poema mediante el hexámetro, las
fórmulas y los neologismos. El capítulo explica cómo su libertad métrica
permite que Homero vuelva a sonar en castellano.

\item \textbf{Helena.}
La Helena de Homero es narradora, hechicera y figura casi divina, capaz de
actuar sin rendir cuentas. Esa impunidad resulta más inquietante que su
conversión cinematográfica en víctima.

\item \textbf{Menelao.}
Un fragmento de \emph{Abandonando Ítaca} ofrece otra versión del matrimonio de
Menelao y Helena. La reinvención del mito sirve para distinguir una lectura
nueva de una simplificación.

\item \textbf{Calipso.}
Los dioses plantean un problema narrativo por su arbitrariedad y su poder para
resolver la historia. Calipso aparece como una mujer sola que retiene a Odiseo
y le ofrece una inmortalidad que él rechaza.

\item \textbf{Nausícaa.}
La princesa feacia reúne juventud, valor y compasión. Su encuentro con Odiseo,
leído desde la adolescencia y desde la paternidad, convierte un amor imposible
en uno de los episodios más humanos del poema.

\item \textbf{Esqueria.}
Los feacios representan el cumplimiento más generoso de la \emph{xenía}.
Odiseo aprende allí a pedir ayuda y se vuelve más humano, razón por la que la
supresión de la isla debilita la tesis de la película.

\item \textbf{Bardos.}
El canto de Demódoco obliga a Odiseo a revelar su identidad. Homero iguala el
llanto del vencedor con el de una viuda troyana y reconoce a la vez el dolor
del verdugo y el de la víctima.

\item \textbf{La Gran Guerra.}
El anciano de \emph{Abandonando Ítaca} recuerda las causas, la duración y el
coste de la guerra de Troya. Su relato contrasta con la versión ennoblecida que
transmiten los bardos.

\item \textbf{Cíclope.}
Polifemo viola conscientemente la hospitalidad y obliga a Odiseo a vencer con
la astucia de «Nadie». La posterior bravuconada del héroe deshace su propia
treta y revela que los dos adversarios comparten inteligencia y arrogancia.

\item \textbf{La bolsa de los vientos.}
La tripulación abre el odre de Eolo cuando Ítaca ya está a la vista y provoca
su ruina. El encuentro con los lestrigones prolonga el tema de la necedad y la
violencia que acaba volviéndose contra quienes la ejercen.

\item \textbf{La isla de la hechicera.}
Circe no es solo una bruja, sino una diosa definida por su voz, su deseo y su
poder. La estancia de un año en Eea sugiere una relación con Odiseo mucho más
compleja que una escala de su viaje.

\item \textbf{La voz de Circe.}
Circe envía a Odiseo al Hades y despide, invisible, a los aqueos tras la
muerte absurda de Elpénor. En los poemas del autor, su voz obliga al héroe a mirarse y abre la
posibilidad de redención.

\item \textbf{Hades.}
Odiseo conversa con Elpénor, Anticlea, Tiresias, Agamenón, Aquiles y Áyax. El
reino de los muertos deshace la gloria heroica y muestra una igualdad final de
la que solo algunos privilegiados escapan.

\item \textbf{Agamenón y Clitemnestra.}
Las versiones de Homero, Esquilo y Eurípides ofrecen culpables y venganzas
distintos. El sacrificio de Ifigenia conduce a una conclusión inequívoca sobre
la imposibilidad de justificar la muerte de un inocente.

\item \textbf{Sirenas.}
Un juicio olímpico absuelve a Nolan por una escena en la que el cine supera a
la palabra. Homero calla qué escucha Odiseo y el poema del autor responde con
un canto sobre la memoria, lo perdido y el tiempo.

\item \textbf{El arco funesto.}
La prueba del arco revela la fuerza de Telémaco y la autonomía de Penélope,
que ayuda al mendigo sin haber reconocido a su marido. El arma anuncia la
matanza sin convertir a Odiseo en un héroe misericordioso.

\item \textbf{La matanza de los pretendientes.}
Odiseo rechaza la reparación ofrecida por Eurímaco y decide exterminar a sus
enemigos. Homero presenta una batalla desigual y feroz que impide confundir al
protagonista con una figura moralmente ejemplar.

\item \textbf{El asesinato de las niñas.}
Las esclavas limpian la sala antes de ser ejecutadas por orden de Odiseo y
Telémaco. El episodio, omitido por el cine, concentra la crueldad del héroe y
la acerca a formas de violencia todavía actuales.

\item \textbf{Penélope.}
El regreso no reúne de inmediato a dos amantes, sino a dos extraños separados
durante veinte años. La prueba del lecho permite el reconocimiento, pero no
borra la ausencia ni todo lo que ambos callan.

\item \textbf{Laertes.}
La visita al padre cierra el regreso familiar de Odiseo y da paso a una
reflexión sobre padres, hijos y memoria. La intervención de Atenea resuelve
finalmente una violencia que los hombres no saben detener.

\item \textbf{Homero.}
La cuestión homérica lleva a la tradición oral, las fórmulas y la memoria de
los bardos. Butler y Graves muestran que los clásicos sobreviven porque cada
época los reescribe, aunque no toda reescritura los comprenda.

\item \textbf{Marcharse de Ítaca.}
Las Ítacas del autor desembocan en un final distinto para el viejo Odiseo. En
lugar de aceptar otra expiación, vuelve al mar para buscar a Circe y afrontar
los fantasmas de su vida.

\end{enumerate}

\end{document}
